% !TeX root = ../papers/nc-p3.tex
\section{異方向二重色族の無印同型の完全判定}
\label{sec:transverse-doublet-isomorphism}
\begin{theorem}\label{thm:transverse-doublet-isomorphism}
定理\ref{thm:klein-sparse-transverse}の全窓について、四元群・対称表示・文字・始点での枠を忘れた
添字保持同型類を次のように分類する。二つの射影パラメータを $p=a/c=(U+3)/(U-1)$,
$q=b/d=(V+3)/(V-1)$ と書き、初めから順序を忘れる。常に $p\ne q$ である。
\begin{enumerate}
\item $p,q$ が有限で $p+q\ne0$ の場合は反対階数16であり、唯一の追加同一視は
$\{p,q\}\sim\{-p,-q\}$ である。
\item 一方が $\infty$ の場合は反対階数12であり、他方の有限パラメータ $q$ の二乗が完全不変量である。
\item $\{p,-p\}$ ($p\ne0$) の場合も反対階数12であるが、前項とは非同型であり、$p^4$ が完全不変量である。
\end{enumerate}
全て最小周期1である。従ってこの族の無印像は既約な二次元の族であり、
射影パラメータの二つの階数12境界での追加同一視も上記で尽くされる。
\end{theorem}

\begin{proof}[可逆な反対切断]
$W$ を前節の対称テンソル $S$ に四語の二コサイクル符号を掛けた表示とする。
$V$ を反対切断 $02\mid13$ の逆行列の転置とすると、四位置の独立な基底変更に対して $V$ は四重双対テンソルとして変換する。
従って $L=W_{01\mid23}V_{01\mid23}^{t}$ は相似変換を受ける。
分子・分母を $c=d=1$ に取り $t=p+q$, $e=p-q$ と置く。
$L$ の特性多項式は
\[
 z^{10}(z-4)^2\big((z-2)^2-A/4\big)^2,\qquad A=(t+4/t)^2 .
\]
$T=\operatorname{tr}_{V_1}(L^2)$ を第1位置について部分traceしたものとすると、
$T$ の二つの対角値はそれぞれ二度現れ、その差は $e(t-4/t)/4$ である。
よって $B=16\big(\operatorname{tr}(T^2)-(\operatorname{tr}T)^2/4\big)=e^2(A-16)$ も無印不変量である。

三つ目は三組の $W,V$ を収縮して
\[
 C=\tfrac12\sum W_{abcd}W_{efgh}W_{ijkl}V^{abgh}V^{ejkd}V^{ifcl}
   =-e^2(1+4/t^2)+\tfrac72(t^2+16/t^2)+100
\]
とする。添字の和は四位置ごとに同じ空間とその双対の間で取る。
\path{computations/verify_transverse_rank_sixteen_isomorphism.py} が特性多項式・部分traceを記号計算する。
最後の収縮では $4tV$ と $W$ の各成分が $p,q$ の一次以下であることを確認した上で、
分母を払った次数6以下の恒等式を各変数7点ずつの整数補間で証明する。

$A\ne16$ なら $e^2=B/(A-16)\ne0$ であり、$t^2+16/t^2=A-8$ と $C$ から $1/t^2$ が一意に決まる。
$A=16$ なら $t^2=4$ であり $C$ が直接 $e^2$ を与える。従って全点で二乗 $t^2,e^2$ が復元される。
$t,e$ の符号の選択は $p,q$ の交換と両符号の同時反転に等しい。
後者は全位置で高位ビットが1の文字を $\sqrt{-1}$ 倍し、全体を $-1$ 倍して実現できる。
\end{proof}

\begin{proof}[二つの特異反対切断の区別]
反対切断の左核を四行四列の行列空間と見て、四次元の核の上に行列式を制限する。
この四次曲面の射影同値類は無印不変である。$p=\infty$ では $a=d=1,c=0,b=q$ として、
核のある座標 $z_0,\ldots,z_3$ で行列式は $-z_2^2 z_3^2$ である。
$\{p,-p\}$ では定数倍を除き
\[
 (z_1^2-z_2z_3)
 \left(p^4(z_2+z_3)^2+4p^2 z_0^2-(z_2-z_3)^2\right)
\]
となる。$p\ne0$ なのでいずれの二次因子も階数3であり、二つの重平面の型に一致しない。
\end{proof}

\begin{proof}[四重交代式とその勾配]
四つの $W$ の各位置をそれぞれ四次交代形式一つに縮約し、第一位置の文字順を固定して正規化した
四次相対不変式を $D$ とする。その勾配を $G$ と書く。
四位置の基底変更に対し $D$ には行列式の積が掛かり、$G/D$ は $D\ne0$ で四重双対テンソルとなる。
上の基底では
\[
 D=16(a^2b^2+a^2d^2+2abcd+b^2c^2+7c^2d^2).
\]
$V=G/D$ として前と同じ $L,T$ を作る。$p=\infty$ の場合に $X=q^2$ と書くと
$D=16(X+1)$ であり、$D=0$ はすでに一つの $X$ を決める。
$X\ne-1$ のとき三つの不変量 $\operatorname{tr}(L^2)/2$, $\operatorname{tr}(L^3)/2$,
$\operatorname{tr}(T^2)-(\operatorname{tr}T)^2/4$ は順に
\[
 \frac{16X^2+40X+11}{8(X+1)^2},\quad
 \frac{64X^3+336X^2+252X+31}{32(X+1)^3},\quad
 \frac{(4X-1)^2}{256(X+1)^4}
\]
となる。三関数について $f(X)-f(Y)$ の分子を取り、
$(X+1)(Y+1)(X-Y)$ を可逆化したイデアルのGröbner基底は $(1)$ である。従って $X$ が復元される。
$\infty$ は符号反転で固定され、有限な $q$ の符号反転は同型となる。
実際、前の高位ビットの倍率で $(a,b,c,d)$ は全体スカラーを除き $(-a,-b,c,d)$ になる。
独立な射影倍率で $[-a:0]=[1:0]$ に戻せる。

$\{p,-p\}$ の場合は $X=p^4$, $D=16(X+7)$ である。$X\ne-7$ のとき $L$ の特性多項式は
\[
 z^{10}\left(z-\frac{X+2}{X+7}\right)^4\left(z-\frac{10}{X+7}\right)^2 .
\]
二つの固有値が合流するのは $X=8$ のみであり、そうでなければ後の固有値が重複度で読める。
前の固有値が零の場合にも後者は非零である。$D=0$ は $X=-7$ だけなので全点で $p^4$ は必要である。

十分性には $p\mapsto\sqrt{-1}\,p$ を実現すればよい。高位ビットを固定して二文字の低位ビットをFourier変換する。
この表示で原始八乗根 $\zeta$ の次の指数を四位置にそれぞれ対角作用させる。
\[
 (-7,-7,-6,-10),\quad(3,3,4,0),\quad(3,3,4,0),\quad(3,3,4,0).
\]
係数が $p$ に比例する非零語では指数和は8を法として2、他の非零語では0である。
従って $p$ だけが $\zeta^2$ 倍になる。
\path{computations/verify_transverse_rank_twelve_isomorphism.py} は二つの核の基底からの行列式、
交代式の勾配、三有理関数の単射性と最後の全係数の位相を検算する。
\end{proof}
